\documentclass[10pt]{article}
\usepackage[margin=20mm]{geometry}
\usepackage{amsmath,amssymb,mathtools,fontspec,unicode-math,xurl,hyperref}
\IfFontExistsTF{DejaVu Serif}{\setmainfont{DejaVu Serif}}{\setmainfont{Latin Modern Roman}}
\setmathfont{latinmodern-math.otf}
\setlength{\parindent}{0pt}\setlength{\parskip}{6pt}\sloppy
\title{Technical review of
The \(2D\) nonlinear shallow water equations
with a partially immersed obstacle}
\author{AI technical verification}
\date{}
\begin{document}
\maketitle
Abstract: This report concerns the 63-page arXiv version 2306.15285v1 by Tatsuo Iguchi and David Lannes. All source pages were read and visually inspected. Three substantive defects are established: the main theorem’s uniform boundary estimate asks for an extra derivative; the compatible-data approximation drops a nonhomogeneous boundary source; and the foundational energy estimate incorrectly retains a final signed flux after taking a time supremum. Fully admissible smooth examples establish each failure. A compatible high-frequency family, an exact stationary solution, and a transported boundary pulse supply the respective evidence. The energy estimate has a directly proved repair. These findings warrant a Major assessment of the presented statements and proof, without asserting that the intended local existence and uniqueness result is false. Two local normalization and indexing corrections are also recorded. Scientific adjudication remains pending.\par
\section*{Source identity and review scope}
The assigned identifier is M25-0263. The frozen title is The \(2D\) nonlinear shallow water equations with a partially immersed obstacle, by Tatsuo Iguchi and David Lannes; the journal DOI is `10.4171/jems/1636`. The actual source is `arXiv:2306.15285v1`, submitted 27 June 2023 at 08:18:03 UTC, with 63 pages. The frozen 2025 batch assignment is retained. No equivalence with the publisher’s version is assumed. All page references below refer to this exact PDF; its physical and printed page numbers agree.\par
The review covers the modeling and potential reduction, definitions of the boundary and time norms, weak dissipativity, generalized vorticity, nonlinear energy estimates, regularization, compatible-data approximation, linear existence, and the final nonlinear closure. Reading every page is distinct from reconstructing every cited external result. The independent derivations below are the decisive checks. Three native AI auxiliaries supplied bounded text and proof checks; the primary reviewer integrated them, independently verified the polar compatibility identities, and performed the recorded source and report image inspections. No human review or independent acceptance is represented.\par
\section*{Assessment and proof dependencies}
The overall assessment is Major. Problem 1 disproves a quantitative conclusion of Theorem 5 under all of its assumptions. Problem 2 invalidates the displayed compatible-approximation construction used in Theorem 4. Problem 3 disproves the basic weak-dissipation estimate from which the higher estimates are derived. Their severity reflects those actual roles; the availability of repairs does not remove the failures as printed.\par
The dependence is particularly concrete. Proposition 6 enters Propositions 7 and 8 and Theorems 1 and 2, then the transformed shallow-water estimates. Proposition 16 supplies Theorem 4, Step 2, and its compatible approximation is invoked again in Section 7.4. Theorem 5 is the final nonlinear statement. The examples do not disprove its intended existence, uniqueness, or the membership in equation (136). Problems 4 and 5 are Minor: the main-body normalization and the actual nonzero-eigenvalue reconstruction already indicate their local corrections.\par
\section*{Confirmed substantive findings}
Problem 1 (Major: an extra boundary derivative in Theorem 5).  On p. 57, after (136), Theorem 5 asserts \[\|u\|_{\mathbb W_T^m}+|\mathfrak d_{\rm tan}\psi_i|_{\mathbb W_{b,T}^{m+1/2}}
 \le C(c_0^{-1},M_0).\] The hypotheses bound the initial tangential derivative only in \(H^{m-1/2}(\Gamma)\). By the definition on p. 15, the displayed conclusion would bound its \(H^{m+1/2}\) norm at time zero. The following construction verifies that compatibility does not supply this missing derivative.\par
\subsection*{A smooth family satisfying every compatibility condition}
Take \(\mathcal I\) to be the unit disk, \(\mathcal E=\{r>1\}\), and \(H_0=h_i=g=1\), so \(\zeta_w=0\) and \(\Lambda=|D_\theta|\). These satisfy Assumption 1. Set \(m=3\) and \(c_0=1/4\). For a fixed sufficiently small \(\delta>0\) and integers \(n\ge2\), let \[\psi_n=\delta n^{-7/2}\cos(n\theta),\qquad
 p=\Lambda\psi_n,\qquad q=\partial_\theta\psi_n.\] Thus \(p'=nq\), \(q'=-np\), and \(p^2+q^2\) is constant. Here primes mean tangential derivatives. Prescribe the initial normal jets at \(r=1\) by \[\phi=\psi_n,\quad \phi_r=p,\quad \phi_{rr}=0,\quad \phi_{rrr}=d,
 \qquad \zeta=0,\quad \zeta_r=a,\quad \zeta_{rr}=0,
 \qquad a=\frac{(n-1)(p^2-q^2)}{1+p^2}.\] Set \(v=\nabla\phi\). The irrotationality and potential matching required for Proposition 9 hold, including tangential matching.\par
Equation reference: R. The initial time jets are generated directly from \[\zeta_t=-\mathop{\mathrm{div}}((1+\zeta)\nabla\phi),\qquad
 \phi_t=-\zeta-\tfrac12|\nabla\phi|^2.\] Write \(G=(1+\zeta)\phi_r-\Lambda(\phi|_\Gamma)\). These recursions are the potential form of (134), and the required conditions (135) are \(G=G_t=G_{tt}=0\). The first two are \[G=p-\Lambda\psi_n=0,\qquad
 G_t=(n-1)(p^2-q^2)-(1+p^2)a=0.\] For completeness, define \[\begin{aligned}
 D&=p+q',& Z&=-D-ap,& U&=-a-qp'+q^2,\\
 Z_{r,0}&=-aD-(-p+p''-2q')-a'q,&
 W&=-3q^2+4qp'-(p')^2,& V&=-Z-pU,\end{aligned}\] and \[
\begin{split}
 R={}&-(1+p^2)Z_{r,0}-2pW-pZD-pqZ'\\
    &+[2Z-p(1+a)]U-qU'-\Lambda V.
\end{split}\] Direct polar differentiation gives \[G_{tt}=(1+3p^2)d+R.\] Choose \(d=-R/(1+3p^2)\). This is a smooth real function and makes the second-order compatibility condition exact.\par
To make this calculation reproducible, the intermediate identities are \[\zeta_{tr}=Z_{r,0}-d,\qquad \phi_{trr}=W-pd,\] \[\phi_{ttr}=(1+p^2)d-Z_{r,0}-pW-qU',\qquad
 \zeta_{tt}=2pd-ZD-pZ_{r,0}-qZ'-W-(1+a)U.\] The trace \(\phi_t=-(p^2+q^2)/2\) is constant in \(\theta\), and \(\phi_{tt}|_\Gamma=V\). Substitution into \(G_{tt}=\zeta_{tt}p+2ZU+\phi_{ttr}-\Lambda V\) gives [R] (\#R). A self-written symbolic differentiation of the polar PDE independently verified these identities with exact zero residuals; it did not execute author code.\par
\subsection*{Uniform norms and the contradiction}
The jets have the uniform bounds \[\|p\|_{H^{5/2}}+\|q\|_{H^{5/2}}\le C\delta,\qquad
 \|a\|_{H^{3/2}}\le C\delta^2,\qquad
 \|d\|_{H^{1/2}}\le C\delta.\] Indeed \(D,Z,U,V\) have the needed \(H^{3/2}\) bounds; \(Z_{r,0}\) and \(R\) have \(H^{1/2}\) bounds. The largest linear second derivative, \(p''\), has amplitude \(\delta n^{-1/2}\) and bounded \(H^{1/2}\) norm. Multiplication by \(H^{3/2}\) functions, the map \(\Lambda:H^{3/2}\to H^{1/2}\), and composition with the positive denominators above preserve these bounds.\par
Here is an explicit extension argument, so that no unspecified compatibility-preserving smoothing is needed. In a fixed collar \(s=r-1\ge0\), choose \(\chi\in C_c^\infty([0,\infty))\) equal to one near zero and define, for prescribed jets \(g_j\), \[E_k(g)(s,\theta)=\sum_{j=0}^{k-1}\frac{s^j}{j!}
                     \chi(s\langle D_\theta\rangle)g_j(\theta).\] All positive derivatives of \(\chi\) vanish at zero, so the prescribed normal jets are exact. Fourier Parseval and the change of variable \(s\mapsto s\langle \ell\rangle\) give \[\|E_k(g)\|_{H^k}\le
 C\sum_{j=0}^{k-1}\|g_j\|_{H^{k-j-1/2}}.\] The polar and flat collar norms are equivalent. Smooth input gives smooth data, supported in the same collar and extending smoothly by zero at its outer edge. Apply \(E_4\) to \((\psi_n,p,0,d)\) and \(E_3\) to \((0,a,0)\). It follows that \[\|\phi_n\|_{H^4(\mathcal E)}+\|\zeta_n\|_{H^3(\mathcal E)}\le C\delta.\] Consequently, for one \(M_0\) independent of \(n\), \[\|(\zeta_n,\nabla\phi_n)\|_{H^3}
       +|\partial_\theta\psi_n|_{H^{5/2}}\le M_0.\] Choosing the fixed \(\delta\) small, Sobolev embedding also gives \(1+\zeta_n-|\nabla\phi_n|^2\ge1/2=2c_0\) throughout \(\mathcal E\). The initial data are therefore smooth, uniformly subcritical, and compatible to orders zero, one and two. Nevertheless \[|\partial_\theta\psi_n|_{H^{7/2}}\asymp\delta n\longrightarrow\infty.\] The printed bound is false already at \(t=0\).\par
\subsection*{Impact and repair boundary}
The excessive uniform estimate is part of the main theorem, rather than a consequence of choosing incompatible data. The proof on pp. 60–61 controls a lower boundary order. A suitable correction uses that order, or explicitly the energy seminorm \[\sup_{0\le t\le T}\sum_{j=0}^{m}
 |\mathfrak d_{\rm tan}\partial_t^j\psi_i(t)|_{H^{m-j-1/2}}.\] The existence and uniqueness conclusion and \(\psi_i\in\mathbb W_{b,T}^{m+1/2}\) are not contradicted. Simply deleting \(\mathfrak d_{\rm tan}\) from the printed bound would introduce another issue: \(M_0\) does not bound the additive constant of \(\psi_i\). Such a full-potential bound needs a gauge or an initial constant-mode bound.\par
Problem 2 (Major: the approximation in Proposition 16 loses \(g_1\)).  Equation (110), p. 48, requires \(N\cdot u_j^{II}=\Lambda\psi_{i,j}
+(\partial_t^j g_1)|_{t=0}\). The correction displayed in the proof of Proposition 16, p. 52, omits this last term. Its following assertion that the uncorrected residual tends to zero is likewise false for general \(g_1\). Theorem 4, Step 2, p. 54, still treats arbitrary \(g_1^{(n)}\); there is no prior homogeneous reduction.\par
Use the unit disk with \(h_i=1\) and the constant matrix \(S=I_3\). Choose a smooth radial \(\chi\) equal to one near \(r=1\) and zero for \(r\ge2\). The time-independent fields and data \[u=(0,\chi(r)e_r),\qquad \psi_i=0,\qquad
 f=(\chi'(r)+\chi(r)/r,0,0),\qquad g_1=1,\quad g_2=0\] solve (103)–(105) exactly: \(G_j\partial_j u=f\), \(N\cdot u^{II}=1=\Lambda\psi_i+g_1\), and \(\partial_t\psi_i+u^I=0\). They have all the stipulated regularity and positivity, and all positive time jets vanish. Thus compatibility holds to every order.\par
Take every approximating datum equal to the original one and \(\widetilde u^{\rm in(n)}=u^{\rm in}\). A valid correction is \(w=0\). The displayed formula on p. 52 instead forces, at order zero, \[N\cdot w_0^{II}
  =-\bigl(N\cdot\widetilde u_0^{II}
                  -\Lambda\widetilde\psi_{i,0}\bigr)=-1.\] The corrected initial normal trace becomes zero and violates the required value \(g_1=1\). Moreover \(w_0\) cannot tend to zero in \(H^m\), because its normal trace is fixed at \(-1\).\par
There is also a sign inconsistency on the same page: \(\phi_i=\widetilde\psi_i-\psi_i\) is defined, whereas the negative-\(w^I\) recursion corresponds to the opposite convention. A consistent correction is to define \(\phi_i=\psi_i-\widetilde\psi_i\) and use \[N\cdot w_j^{II}=\Lambda\phi_{i,j}
 +g_{1,j}^{(n)}
 -\bigl(N\cdot\widetilde u_j^{II}
                     -\Lambda\widetilde\psi_{i,j}\bigr).\] The example then requires zero correction as expected. All subsequent jets and extension steps must use this convention.\par
This disproves the displayed construction, not the existential statement of Proposition 16. Its role in constructing compatible smooth data for Theorem 4 makes the gap substantive. Nonzero commutator boundary data also occur in the final tangential application on p. 61. A complete corrected approximation proof has not been certified by this review.\par
Problem 3 (Major: a final signed flux does not bound the past peak).  Equation (45), p. 21, and Proposition 6, p. 22, use the final cumulative signed boundary flux to control the norm \(I_{\lambda,t}\), which includes a supremum over earlier times. Outflow can cancel earlier inflow without erasing that supremum. The following smooth example satisfies Assumption 2 and Definition 2.\par
Let \(\mathcal E=\{r>1\}\), use a scalar unknown, and set \(S=1\), \(B=f=0\). Put \[T_*=\frac1{100},\qquad b=\frac{T_*}{8\pi},\qquad
 s(t)=-2b\cos(2\pi t/T_*).\] Then \[a(t)=s'(t)=\tfrac12\sin(2\pi t/T_*),
 \qquad A_j=a(t)\frac{x_j}{r}.\] All coefficients are smooth. Since \(|a|\le1/2\) and \(\mathop{\mathrm{div}}A=a/r\le1/2\), Assumption 2 holds with \(\alpha_0=\beta_0=\beta_0^{\rm in}=\beta_1=1\). That assumption does not bound \(\partial_t A\). For nonzero \(\eta\in C_c^\infty(-b,b)\), let \[u(t,x)=\eta(r-1-s(t)).\] Then \(u_t+A_j\partial_j u=0\), \(u(0)=u(T_*)=0\), and \(u\) is nonzero at intermediate times. Its support is in a fixed boundary collar. Prescribe its smooth boundary trace \(g_b=u|_\Gamma\).\par
For every \(\lambda\ge1\), define \[S_{\rm data}(t)=\int_0^t e^{-2\lambda\tau}
                  \int_\Gamma a_+|g_b|^2,\qquad
 E_{\rm bdry}(t)=\int_0^t e^{-2\lambda\tau}
                  \int_\Gamma a_-|u|^2 .\] These are nonnegative; \(S_{\rm data}\) is nondecreasing and depends only on the prescribed data and coefficients. The boundary quadratic form is \(\mathcal B[u]=a u^2\), and its integral equals \(S_{\rm data}-E_{\rm bdry}\). Definition 2 therefore holds even without its permitted bulk remainder. The \(\lambda\) dependence is of the same kind as in the paper’s Examples 2 and 3.\par
Take \(\lambda=1\), and write \(E(t)=\|u(t)\|_2^2\) and \(M=\max_{[0,T_*]}e^{-2t}E(t)>0\). The exterior outward normal is \(-e_r\), so integration by parts gives \[E'(t)=\int_\Gamma a u^2+\int_\mathcal E\frac{a(t)}r u^2.\] Since both endpoint energies vanish, \[F:=\int_0^{T_*}e^{-2t}\int_\Gamma a u^2
 =\int_0^{T_*}e^{-2t}\int_\mathcal E(2-a/r)u^2
 \le\tfrac52 T_*M.\] Definition (28) yields \(I_{1,T_*}(\|u\|_2)^2\ge M\). Equation (45), with zero initial and body-source data, would give \(I^2\le4F\le M/10\). Proposition 6 would give \(I^2\le8(S_{\rm data}-E_{\rm bdry})=8F\le M/5\). Both are contradictions. No choice of the unspecified constant on the vanishing initial and body-source terms can fix them.\par
\subsection*{A proved replacement estimate}
Equation reference: repair. One must insert weak dissipativity into the endpoint energy identity before taking a supremum. Here is a sufficient repair with explicit, nonoptimal constants. Write \[v(t)=(S(t)u(t),u(t))^{1/2},\quad J=I_{\lambda,t}(v),\quad
 A=(S(0)u(0),u(0)),\quad
 B=\sqrt{\beta_0}\,S^*_{\lambda,t}(\|f\|_2).\] For \(\lambda\ge\beta_1/\alpha_0\), the endpoint inequality and Definition 2 bound, separately, \[X=\sup_{\tau\le t}e^{-2\lambda\tau}v(\tau)^2,\qquad
 Y+E_{\rm bdry}(t)
       =\lambda\int_0^t e^{-2\lambda\tau}v(\tau)^2\,d\tau
          +E_{\rm bdry}(t)\] by \(Q=A+2BJ+S_{\rm data}(t)+J^2/8\). This uses nonnegativity of the boundary energy and monotonicity of the data and source norms; it does not require monotonicity of \(E_{\rm bdry}\). Since \(J^2\le2(X+Y)\), it follows that \[\tfrac12J^2+2E_{\rm bdry}\le4A+8BJ+4S_{\rm data}.\] Young’s inequality \(8BJ\le J^2/4+64B^2\) then gives \[
\begin{split}
 I_{\lambda,t}(\|u\|_2)^2+\frac8{\alpha_0}E_{\rm bdry}
 \le{}&16\frac{\beta_0^{\rm in}}{\alpha_0}\|u(0)\|_2^2\\
 &+256\frac{\beta_0}{\alpha_0}
             \bigl(S^*_{\lambda,t}(\|f\|_2)\bigr)^2
 +\frac{16}{\alpha_0}S_{\rm data}.
\end{split}\] There is no added derivative or coefficient regularity requirement. The foundational estimate and its proof need correction, but [repair] (\#repair) preserves the qualitative energy control. Later estimates with an unspecified data constant can accommodate such a change. This is not a counterexample to nonlinear local well-posedness; full propagation through every later constant is not asserted.\par
\section*{Confirmed local corrections}
Problem 4 (Minor: pressure and fixed-energy normalization).  The Bernoulli row in introductory (3), p. 4, omits the density factor present in the physical momentum equations on p. 3 and the consistent main-body formulation on p. 8. Take the unit disk, \(H_0=2\), \(h_i=1\), \(\zeta_w=-1\), \(g=1\), \(\rho=2\), \(P_{\rm atm}=0\), \(\zeta=v=v_i=\phi=\psi_i=0\), and \(\underline P_i=2\). These smooth stationary data satisfy the equations, positive depths, irrotationality, all matching and compatibility conditions, and strict subcriticality. The body uses \(\underline P_i+\rho g\zeta_w=0\), which holds. Equation (3) instead requires \(0=\underline P_i+g\zeta_w=1\). Divide its pressure terms by \(\rho\), or multiply its gravitational and kinetic terms by \(\rho\). The main-body argument already uses the consistent convention.\par
On p. 31, the equality of reduced energy with \(\mathfrak E_{\rm fluid}/\rho\) also misses the fixed term \((g/2)\int_\mathcal I\zeta_w^2\). For the same equilibrium the reduced energy is zero whereas \(\mathfrak E_{\rm fluid}/\rho=\pi/2\). State equality up to this constant, or specify a relative-energy normalization. Its time derivative is zero, so the conservation argument is unaffected.\par
Problem 5 (Minor: the first inverse-eigenvalue index).  Assumption 6(i), p. 27, starts its inverse-eigenvalue list at \(\lambda_{n_1}^{-1}\). The intended nonzero block starts at \(n_1+1\), as in Lemma 3. Corollary 2 on p. 39 uses \(n_1=1\), \(n_2=2\), and eigenvalues \(0,1,-1\); these occur already for the constant positive matrix \(S=I_3\). The literal list would therefore require the reciprocal of zero in the very application it is meant to cover. Change its first index to \(n_1+1\). This is a mismatch between a written assumption and its application, not a counterexample to a theorem conditional on a stronger assumption. The actual reconstruction only divides by nonzero eigenvalues.\par
\section*{Checks without an established defect}
The following checks delimit the findings. The interface vorticity flux cancellation follows from mass-flux matching and does not divide by a potentially vanishing normal velocity. The single Jordan boundary, zero interior vorticity, and tangential matching give zero exterior circulation, so the potential reduction includes the needed topological condition. Green’s identity gives positivity and self-adjointness of \(\Lambda\); its constant kernel is correctly handled by a tangential derivative seminorm.\par
The transformation by \(\Sigma(u)\), the constant matrices \(G_j\), and \(S=\Sigma^{-1}\) was checked algebraically. The generalized vorticity identities cancel the second spatial derivatives; the normal reconstruction needs the invertibility of \(L\), not an extra assumption that \(\widetilde A_0\) is invertible. Applying Lemma 3 to differentiated unknowns retains the original \(Q,W\), so this step does not itself demand an unlisted high derivative of \(Q\). The two-dimensional product estimates and the use of order \(m-1\) in the nonlinear application give the stated derivative count in the checked terms.\par
For \(m=3\), a uniform \(\mathbb W_T^3\) bound and strong \(\mathbb W_T^2\) continuity yield \(\mathbb W_T^{2,p}\) continuity for every finite \(p\) by interpolation. This supplies the weaker coefficient regularity used in Theorem 4. The two-dimensional acoustic normal block has determinant \((N\cdot v)^2-gh<0\) under subcriticality; conserved curl controls the remaining component. In Proposition 17, differentiated bulk and boundary residuals vanish, and (139) sets their initial values to zero. No separate failure of this recovery step was found.\par
\section*{Coverage, limitations and disposition}
All 63 source pages, including the bibliography, were read and viewed. The checks above reconstruct selected decisive arguments; they are not a new proof of all the external hyperbolic existence, elliptic regularity and Sobolev results cited by the authors. The high-frequency example proves a false uniform estimate, the stationary example proves a failed construction, and the pulse example proves a false energy inequality. None is being enlarged into an unproved assertion about every formulation or version of the main theorem.\par
Bounded searches for supplements, code and corrections produced no applicable separate item. This is a search limitation, not proof of absence and not a mathematical defect. The publisher text was not obtained for a version comparison. An exploratory concern about Proposition 5’s parameter dependence and a possible Proposition 15 sign discrepancy were not established under all hypotheses and are excluded from the findings.\par
Prior false or incomplete checkpoints remain historical records. Author adoption is unknown. The initial scientific decision remains pending; the human-review list is empty.\par
9 T. Iguchi and D. Lannes, The \(2D\) nonlinear shallow water equations with a partially immersed obstacle, arXiv:2306.15285v1, 27 June 2023, 63 pp. Journal identifier: DOI `10.4171/jems/1636`. The exact version reviewed here is the arXiv manuscript.\par
\end{document}
